\section{The Lie Algebra of Rotations}

The quantization of angular momentum is structurally defined by the
non-commutative geometry of the rotation group $\mathrm{SO}(3)$. We derive the Lie
algebra $\mathfrak{so}(3)$ directly from the canonical commutation relations
(CCR) of position and momentum.

\subsection{Fundamental Commutation Relations}

We compute the commutator $[\hat{L}_x, \hat{L}_y]$ using the operator definition
$\hat{L}_i = \epsilon_{ijk} \hat{x}_j \hat{p}_k$ and the Heisenberg algebra
$[\hat{x}_i, \hat{p}_j] = i\hbar \delta_{ij}$.

\subsubsection{Algebraic Derivation}
Expanding the components $\hat{L}_x = \hat{y}\hat{p}_z - \hat{z}\hat{p}_y$ and
$\hat{L}_y = \hat{z}\hat{p}_x - \hat{x}\hat{p}_z$:
\begin{equation}
  [\hat{L}_x, \hat{L}_y] = [\hat{y}\hat{p}_z - \hat{z}\hat{p}_y,
  \hat{z}\hat{p}_x - \hat{x}\hat{p}_z].
\end{equation}
Using the linearity of the commutator and the product rule $[AB, C] = A[B,C] +
[A,C]B$, we identify the non-vanishing terms. The only non-commuting pairs
involve coordinates and momenta of the same index (specifically $z$ and $p_z$
here).
\begin{align}
  [\hat{L}_x, \hat{L}_y] &= [\hat{y}\hat{p}_z, \hat{z}\hat{p}_x] +
  [-\hat{z}\hat{p}_y, -\hat{x}\hat{p}_z] \quad (\text{Cross terms vanish})
  \nonumber \\
                         &= \hat{y}[\hat{p}_z, \hat{z}]\hat{p}_x + \hat{x}\hat{p}_y[\hat{z},
                         \hat{p}_z] \nonumber \\
                         &= \hat{y}(-i\hbar)\hat{p}_x + \hat{x}\hat{p}_y(i\hbar) \nonumber \\
                         &= i\hbar (\hat{x}\hat{p}_y - \hat{y}\hat{p}_x) = i\hbar \hat{L}_z.
\end{align}
Cyclic permutation of indices $(x \to y \to z)$ yields the complete Lie algebra
of the group:
\begin{equation} \label{eq:lie_algebra}
  [\hat{L}_i, \hat{L}_j] = i\hbar \epsilon_{ijk} \hat{L}_k.
\end{equation}

\subsubsection{Geometric Interpretation}

The non-commutativity $[\hat{L}_x, \hat{L}_y] \neq 0$ reflects the non-Abelian
nature of 3D rotations: a rotation about the x-axis followed by the y-axis is
distinct from the reverse order. In the language of Geometric Algebra, this
arises because the generators are bivectors (oriented planes); the plane
$xy$ (generated by $L_z$) is not orthogonal to the interaction of planes $yz$
and $zx$.

\subsection{Commutation of the Casimir Invariant \texorpdfstring{$\hat{L}^2$}{L²}}

The total angular momentum operator is defined as the contraction of the vector
operator with itself:
\begin{equation}
  \hat{L}^2 \equiv \hat{\mathbf{L}} \cdot \hat{\mathbf{L}} = \sum_{i}
  \hat{L}_i^2.
\end{equation}
This operator is a scalar (rank-0 tensor) and represents the Casimir invariant
of the $\mathfrak{so}(3)$ algebra.

\subsubsection{Commutation with Generators}
To verify that $\hat{L}^2$ commutes with the generators, we compute $[\hat{L}^2,
\hat{L}_k]$. Using the identity $[\hat{A}^2, \hat{B}] = \hat{A}[\hat{A}, \hat{B}]
+ [\hat{A}, \hat{B}]\hat{A}$:
\begin{align}
  [\hat{L}^2, \hat{L}_z] &= [\hat{L}_x^2 + \hat{L}_y^2 + \hat{L}_z^2,
  \hat{L}_z] \nonumber \\
                         &= [\hat{L}_x^2, \hat{L}_z] + [\hat{L}_y^2, \hat{L}_z] \quad (\text{since }
                         [\hat{L}_z, \hat{L}_z]=0) \nonumber \\
                         &= \hat{L}_x[\hat{L}_x, \hat{L}_z] + [\hat{L}_x, \hat{L}_z]\hat{L}_x +
                         \hat{L}_y[\hat{L}_y, \hat{L}_z] + [\hat{L}_y, \hat{L}_z]\hat{L}_y.
\end{align}
Substituting the structure constants from Eq. (\ref{eq:lie_algebra}):
\begin{itemize}
  \item $[\hat{L}_x, \hat{L}_z] = -i\hbar \hat{L}_y$
  \item $[\hat{L}_y, \hat{L}_z] = +i\hbar \hat{L}_x$
\end{itemize}
\begin{align}
  [\hat{L}^2, \hat{L}_z] &= -i\hbar(\hat{L}_x \hat{L}_y + \hat{L}_y \hat{L}_x)
  + i\hbar(\hat{L}_y \hat{L}_x + \hat{L}_x \hat{L}_y) \nonumber \\
                         &= 0.
\end{align}
By symmetry, $[\hat{L}^2, \hat{\mathbf{L}}] = 0$.

\subsubsection{Physical Significance: Compatible Observables}
The commutation $[\hat{L}^2, \hat{L}_z] = 0$ implies that the magnitude of
rotation (related to the radius of the sphere) and the plane of rotation
(orientation) are compatible physical properties.
\begin{itemize}
  \item Scalar Nature: $\hat{L}^2$ is a scalar operator; it measures the
    ``amount'' of rotation invariant under coordinate changes.
  \item Vector Nature: $\hat{L}_z$ is a component of a vector (or
    bivector) and depends on the choice of axis.
\end{itemize}
Because they commute, they share a common basis of eigenstates $\ket{\ell, m}$,
allowing for the simultaneous quantization of total energy, total angular
momentum, and azimuthal orientation.
